% TOP_PROBLEM: {"id": 7500110, "title": "Strength of the Dual Ramsey Theorem", "category_id": 18, "category": {"id": 18, "name": "logic", "display_name": "Logic", "description": "Problems in mathematical logic, model theory, proof theory, and finite model theory.", "slug": "logic", "order_index": 18, "created_at": "2026-07-31T15:26:25.671Z"}, "status": "partially_solved"}
% FIRST_POSTED: null
% ENTRY_KIND: statement_only
% Imported from: batch20.tex; UnsolvedMath ID 7500110
\documentclass[11pt]{article}
\usepackage[a4paper,margin=25mm]{geometry}
\usepackage{amsmath,amssymb,mathtools}
\usepackage{fontspec,unicode-math}
\setmainfont{Latin Modern Roman}
\setsansfont{Latin Modern Sans}
\setmathfont{Latin Modern Math}
\usepackage{xurl,hyperref}
\hypersetup{colorlinks=true,linkcolor=blue,urlcolor=blue,pdftitle={UnsolvedMath Problem 7500110}}
\newcommand{\sref}[2]{\hyperlink{source:#1:#2}{[S#2]}}
\newcommand{\cataloguescope}[1]{\paragraph{Mathematical question.}#1}
\begin{document}
% UnsolvedMath ID 7500110; source catalogue code AMR-074-0110
\setcounter{section}{7500109}
\section{Strength of the Dual Ramsey Theorem}\label{7500110}
{\small\sffamily UnsolvedMath ID 7500110 · Logic}
\subsection{Definitions and mathematical statement}

\paragraph{Shared notation.}$\mathbb N=\{1,2,\ldots\}$, and $\mathbb Z,\mathbb Q,\mathbb R,\mathbb C$ denote the integers, rationals, reals and complex numbers. Any other conventions are specified locally.


Work in second-order arithmetic over $\mathsf{RCA}_0$, the base theory with $\Delta^0_1$ comprehension and $\Sigma^0_1$ induction. The strength of a principle means which subsystems prove it and which subsystems it implies over that base; equivalence requires both implications. For a positive finite $k$ or $k=\omega$, $(\omega)^k$ is the space of partitions of $\mathbb N$ into $k$ nonempty pieces, with its usual product coding by assigning each integer its block label, labels ordered by the least member of each block. A Borel colouring is a finite-valued colouring whose colour classes are Borel in this coding. If $X\in(\omega)^\omega$, let $(X)^k$ be the partitions into $k$ pieces coarser than $X$, so every piece is a union of pieces of $X$. The dual Ramsey principle $\mathsf{DRT}_k$ says that every finite Borel colouring of $(\omega)^k$ has an $X\in(\omega)^\omega$ for which $(X)^k$ is monochromatic.
\paragraph{Statement recorded in the supplied source.}
Determine the reverse-mathematical strength of $\mathsf{DRT}_k$, for $k\leq\omega$. \sref{7500110}{2}
\subsection{Short English statement}
Determine which set-existence axioms are required for the Borel dual Ramsey theorem for partitions into finitely many, or countably many, pieces.
\subsection{Sources}
\begin{description}
\item[\hypertarget{source:7500110:1}{S1}] ULAM.AI and the UnsolvedMath Contributors, UnsolvedMath: A Curated Collection of Open Mathematics Problems, record 7500110 (AMR-074-0110). CC BY 4.0. \url{https://huggingface.co/datasets/ulamai/UnsolvedMath}.
\item[\hypertarget{source:7500110:2}{S2}] Antonio Montalbán, Open Questions in Reverse Mathematics, author preprint dated 17 August 2010, Question 10, printed p. 5. \url{https://math.berkeley.edu/~antonio/papers/questionsRM.pdf}.
\end{description}
\label{end:7500110}
\end{document}
